% Generated from lessons/0012-2-chain-rule-multivariable.html; edit the HTML source, then rerun the converter.
\section{多元复合函数的求导法则}
\lessonmark{多元复合函数的求导法则}

这一节把一元链式法则推广到多元函数。考试中最重要的不是背公式，而是画清楚“外层函数依赖哪些中间变量，中间变量又依赖哪些自变量”，再按路径把导数相乘相加。

\subsection*{本节主线}

\begin{conceptbox}
\textbf{链式法则}
沿每条依赖路径求导，把所有路径贡献相加。
\end{conceptbox}

\begin{conceptbox}
\textbf{Jacobi 矩阵乘法}
向量值函数复合时，导数矩阵按复合顺序相乘。
\end{conceptbox}

\begin{conceptbox}
\textbf{微分形式}
一阶全微分形式不变；高阶微分在变量替换时要格外小心。
\end{conceptbox}

\begin{notebox}
一句话：复合函数求导就是“顺着依赖图走路径”。一阶问题看路径相加；向量值问题看矩阵相乘；二阶问题多考虑一层“导数本身也在复合”。
\end{notebox}

\subsection*{一、二元到二元的链式法则}

设

\[
      z=f(x,y),\qquad x=x(u,v),\quad y=y(u,v).
    \]

于是 \(z\) 也成为 \(u,v\) 的函数：

\[
      z=f(x(u,v),y(u,v)).
    \]

\subsubsection*{定理 12.2.1：链式法则}

若 \(x=x(u,v),y=y(u,v)\) 在 \((u_0,v_0)\) 可导，且 \(f\) 在 \((x_0,y_0)\) 可微，其中

\[
      x_0=x(u_0,v_0),\qquad y_0=y(u_0,v_0),
    \]

则复合函数 \(z=f(x(u,v),y(u,v))\) 对 \(u,v\) 可偏导，且

\[
      \frac{\partial z}{\partial u}
      =
      \frac{\partial f}{\partial x}\frac{\partial x}{\partial u}
      +
      \frac{\partial f}{\partial y}\frac{\partial y}{\partial u},
      \qquad
      \frac{\partial z}{\partial v}
      =
      \frac{\partial f}{\partial x}\frac{\partial x}{\partial v}
      +
      \frac{\partial f}{\partial y}\frac{\partial y}{\partial v}.
    \]

右边的 \(\frac{\partial f}{\partial x},\frac{\partial f}{\partial y}\) 要在中间点 \((x(u,v),y(u,v))\) 处取值。

\subsubsection*{证明思路}

因为 \(f\) 可微，有

\[
      \Delta z=
      \frac{\partial f}{\partial x}\Delta x
      +
      \frac{\partial f}{\partial y}\Delta y
      +o\!\left(\sqrt{(\Delta x)^2+(\Delta y)^2}\right).
    \]

再把 \(\Delta x,\Delta y\) 分别看成由 \(\Delta u\) 或 \(\Delta v\) 引起的增量，除以 \(\Delta u\) 或 \(\Delta v\)，令其趋于 \(0\)，就得到公式。

\begin{notebox}
教材特别提醒：这里要求外层 \(f\) 在对应点可微，不能只把条件改成“可偏导”。上一节的反例说明，偏导存在不足以支撑链式法则。
\end{notebox}

\subsection*{二、一般链式法则与矩阵形式}

设

\[
      z=f(y_1,\ldots,y_m),\qquad
      y_j=y_j(x_1,\ldots,x_n).
    \]

若外层 \(f\) 在 \(y=g(x)\) 处可微，内层 \(g\) 在 \(x\) 处可导，则

\[
      \frac{\partial z}{\partial x_i}
      =
      \sum_{j=1}^{m}
      \frac{\partial f}{\partial y_j}
      \frac{\partial y_j}{\partial x_i},
      \qquad i=1,\ldots,n.
    \]

写成矩阵就是

\[
      \begin{pmatrix}
      \frac{\partial z}{\partial x_1}&\cdots&\frac{\partial z}{\partial x_n}
      \end{pmatrix}
      =
      \begin{pmatrix}
      \frac{\partial f}{\partial y_1}&\cdots&\frac{\partial f}{\partial y_m}
      \end{pmatrix}
      \begin{pmatrix}
      \frac{\partial y_1}{\partial x_1}&\cdots&\frac{\partial y_1}{\partial x_n}\\
      \vdots&\ddots&\vdots\\
      \frac{\partial y_m}{\partial x_1}&\cdots&\frac{\partial y_m}{\partial x_n}
      \end{pmatrix}.
    \]

如果 \(f:R^k\to R^m\)，\(g:R^n\to R^k\)，则向量值函数的链式法则是

\[
      J_{f\circ g}(x)=J_f(g(x))\,J_g(x).
    \]

\begin{notebox}
矩阵顺序不能反：先由 \(x\) 进入 \(g\)，再进入 \(f\)，所以导数矩阵是外层在左，内层在右。
\end{notebox}

\subsection*{三、考试算法：先画依赖图}

例如

\[
      w=f(x,u,v),\qquad u=g(y,z),\qquad v=h(x,y).
    \]

求 \(\frac{\partial w}{\partial x}\) 时，有两条路径：

\begin{itemize}
 \item 直接路径：\(x\to w\)
 \item 间接路径：\(x\to v\to w\)
 \end{itemize}

因为 \(u=g(y,z)\) 不含 \(x\)，所以 \(x\to u\to w\) 这条路没有贡献。

\[
      \frac{\partial w}{\partial x}
      =
      \frac{\partial f}{\partial x}
      +
      \frac{\partial f}{\partial v}
      \frac{\partial h}{\partial x}.
    \]

同理

\[
      \frac{\partial w}{\partial y}
      =
      \frac{\partial f}{\partial u}
      \frac{\partial g}{\partial y}
      +
      \frac{\partial f}{\partial v}
      \frac{\partial h}{\partial y},
      \qquad
      \frac{\partial w}{\partial z}
      =
      \frac{\partial f}{\partial u}
      \frac{\partial g}{\partial z}.
    \]

\begin{notebox}
做题时不要急着写公式。先写“外层变量”和“中间变量”，再按路径相乘相加。
\end{notebox}

\subsection*{四、一阶全微分的形式不变性}

若 \(z=f(x,y)\)，当 \(x,y\) 是自变量时，

\[
      dz=
      \frac{\partial z}{\partial x}dx+
      \frac{\partial z}{\partial y}dy.
    \]

若 \(x=x(u,v),y=y(u,v)\) 是中间变量，由链式法则仍有

\[
      dz=
      \frac{\partial z}{\partial x}dx+
      \frac{\partial z}{\partial y}dy.
    \]

这叫一阶全微分形式不变性。它允许我们在复杂表达式里先取对数、先设中间变量，再统一微分。

\subsubsection*{高阶微分不再简单形式不变}

如果 \(x,y\) 是自变量，二阶全微分是

\[
      d^2z=
      \frac{\partial^2z}{\partial x^2}dx^2
      +2\frac{\partial^2z}{\partial x\partial y}dxdy
      +\frac{\partial^2z}{\partial y^2}dy^2.
    \]

但若 \(x,y\) 是中间变量，还要多出

\[
      \frac{\partial z}{\partial x}d^2x+
      \frac{\partial z}{\partial y}d^2y.
    \]

\begin{notebox}
所以一阶全微分可以放心换变量；二阶及以上要检查中间变量的二阶微分是否为零。
\end{notebox}

\subsection*{五、即时判断练习}

\begin{quickcheck}
链式法则中，外层函数只要求偏导存在就够了吗？

\tcblower
\textbf{答案：}外层可微
\end{quickcheck}

\begin{quickcheck}
向量值函数复合 \(f\circ g\) 的 Jacobi 矩阵乘法顺序是什么？

\tcblower
\textbf{答案：}外左内右
\end{quickcheck}

\begin{quickcheck}
全微分形式不变性对高阶微分也直接成立吗？

\tcblower
\textbf{答案：}一阶成立
\end{quickcheck}

\subsection*{六、课后题训练}

下面按教材第 144-146 页习题 1-18 整理。题面完整保留，提示只给入口；写作业时建议先独立算一遍，再展开提示核对方向。

\begin{exercisebox}
\textbf{习题 1\badge{链式法则求偏导}}

利用链式法则求偏导数：

\begin{enumerate}
 \item \(z=\tan(3t+2x^2-y^2),\ x=\frac1t,\ y=\sqrt t\)，求 \(\frac{dz}{dt}\)；
 \item \(z=e^{x-2y},\ x=\sin t,\ y=t^3\)，求 \(\frac{d^2z}{dt^2}\)；
 \item \(w=\frac{e^{ax}(y-z)}{a^2+1},\ y=a\sin x,\ z=\cos x\)，求 \(\frac{dw}{dx}\)；
 \item \(z=u^2\ln v,\ u=\frac{x}{y},\ v=3x-2y\)，求 \(\frac{\partial z}{\partial x},\frac{\partial z}{\partial y}\)；
 \item \(u=e^{x^2+y^2+z^2},\ z=y^2\sin x\)，求 \(\frac{\partial u}{\partial x},\frac{\partial u}{\partial y}\)；
 \item \(w=(x+y+z)\sin(x^2+y^2+z^2),\ x=te^s,\ y=e^t,\ z=e^{s+t}\)，求 \(\frac{\partial w}{\partial s},\frac{\partial w}{\partial t}\)；
 \item \(z=x^2+y^2+\cos(x+y),\ x=u+v,\ y=\arcsin v\)，求 \(\frac{\partial z}{\partial u},\frac{\partial^2 z}{\partial v\partial u}\)；
 \item 以下假设 \(f\) 具有二阶连续偏导数。若 \(u=f\!\left(xy,\frac{x}{y}\right)\)，求 \(\frac{\partial u}{\partial x},\frac{\partial u}{\partial y},\frac{\partial^2u}{\partial x\partial y},\frac{\partial^2u}{\partial y^2}\)；
 \item 若 \(u=f(x^2+y^2+z^2)\)，求 \(\frac{\partial u}{\partial x},\frac{\partial u}{\partial y},\frac{\partial u}{\partial z},\frac{\partial^2u}{\partial x^2},\frac{\partial^2u}{\partial x\partial y}\)；
 \item 若 \(w=f(x,y,z),\ x=u+v,\ y=u-v,\ z=uv\)，求 \(\frac{\partial w}{\partial u},\frac{\partial w}{\partial v},\frac{\partial^2w}{\partial u\partial v}\)。
 \end{enumerate}

\begin{hintbox}
先列出中间变量。比如第 (8) 题设 \(s=xy,\ t=x/y\)，再把 \(u=f(s,t)\) 求导。二阶题要继续对一阶结果求导，注意 \(s,t\) 仍然是 \(x,y\) 的函数。
\end{hintbox}
\end{exercisebox}

\begin{exercisebox}
\textbf{习题 2\badge{沿曲线恒等式}}

设 \(f(x,y)\) 具有连续偏导数，且 \(f(x,x^2)=1,\ \frac{\partial f}{\partial x}(x,x^2)=x\)，求 \(\frac{\partial f}{\partial y}(x,x^2)\)。

\begin{hintbox}
对恒等式 \(f(x,x^2)=1\) 两边关于 \(x\) 求导。
\end{hintbox}
\end{exercisebox}

\begin{exercisebox}
\textbf{习题 3\badge{函数嵌函数}}

设 \(f(x,y)\) 具有连续偏导数，且 \(f(1,1)=1,\ \frac{\partial f}{\partial x}(1,1)=2,\ \frac{\partial f}{\partial y}(1,1)=3\)。如果 \(\varphi(x)=f(x,f(x,x))\)，求 \(\varphi'(1)\)。

\begin{hintbox}
先设 \(g(x)=f(x,x)\)，则 \(\varphi(x)=f(x,g(x))\)。先求 \(g'(1)\)，再求 \(\varphi'(1)\)。
\end{hintbox}
\end{exercisebox}

\begin{exercisebox}
\textbf{习题 4\badge{特殊组合消去}}

设 \(z=\dfrac{y}{f(x^2-y^2)}\)，其中 \(f(t)\) 具有连续导数，且 \(f(t)\ne0\)，求

\[
        \frac1x\frac{\partial z}{\partial x}
        +
        \frac1y\frac{\partial z}{\partial y}.
      \]

\begin{hintbox}
令 \(s=x^2-y^2\)。先把 \(z=y/f(s)\) 看清楚，分别求 \(\frac{\partial z}{\partial x},\frac{\partial z}{\partial y}\)，再代入组合式，很多项会抵消。
\end{hintbox}
\end{exercisebox}

\begin{exercisebox}
\textbf{习题 5\badge{验证恒等式}}

设 \(z=\arctan\frac{x}{y},\ x=u+v,\ y=u-v\)，验证：

\[
        \frac{\partial z}{\partial u}
        +
        \frac{\partial z}{\partial v}
        =
        \frac{u-v}{u^2+v^2}.
      \]

\begin{hintbox}
既可以先把 \(z\) 写成 \(u,v\) 的显式函数再求导，也可以先求 \(\frac{\partial z}{\partial x},\frac{\partial z}{\partial y}\)，再用链式法则。
\end{hintbox}
\end{exercisebox}

\begin{exercisebox}
\textbf{习题 6\badge{偏微分方程验证}}

设 \(\varphi\) 和 \(\psi\) 具有二阶连续导数，验证：

\begin{enumerate}
 \item \(u=y\varphi(x^2-y^2)\) 满足 \(y\frac{\partial u}{\partial x}+x\frac{\partial u}{\partial y}=\frac{x}{y}u\)；
 \item \(u=\varphi(x-at)+\psi(x+at)\) 满足波动方程 \(\frac{\partial^2u}{\partial t^2}=a^2\frac{\partial^2u}{\partial x^2}\)。
 \end{enumerate}

\begin{hintbox}
第 (1) 题令 \(s=x^2-y^2\)。第 (2) 题令 \(\xi=x-at,\eta=x+at\)，分别求二阶偏导。
\end{hintbox}
\end{exercisebox}

\begin{exercisebox}
\textbf{习题 7\badge{Laplacian 变量替换}}

设 \(z=f(x,y)\) 具有二阶连续偏导数，写出

\[
        \frac{\partial^2 z}{\partial x^2}
        +
        \frac{\partial^2 z}{\partial y^2}
      \]

在坐标变换

\[
        u=x^2-y^2,\qquad v=2xy
      \]

下的表达式。

\begin{hintbox}
把 \(z\) 看作 \(z=f(u,v)\)。先写 \(\frac{\partial z}{\partial x},\frac{\partial z}{\partial y}\)，再分别求二阶。混合项会抵消，最后出现 \(4(x^2+y^2)\)。
\end{hintbox}
\end{exercisebox}

\begin{exercisebox}
\textbf{习题 8\badge{积分上限复合}}

设

\[
        f(x,y)=\int_0^{xy}e^{-t^2}\,dt,
      \]

求

\[
        \frac{x}{y}\frac{\partial^2 f}{\partial x^2}
        -2\frac{\partial^2 f}{\partial x\partial y}
        +\frac{y}{x}\frac{\partial^2 f}{\partial y^2}.
      \]

\begin{hintbox}
先设 \(F(s)=\int_0^s e^{-t^2}\,dt\)，则 \(f(x,y)=F(xy)\)，且 \(F'(s)=e^{-s^2}\)。
\end{hintbox}
\end{exercisebox}

\begin{exercisebox}
\textbf{习题 9\badge{齐次函数 Euler 公式}}

如果函数 \(f(x,y)\) 满足：对于任意实数 \(t\) 及 \(x,y\)，成立

\[
        f(tx,ty)=t^n f(x,y),
      \]

那么 \(f\) 称为 \(n\) 次齐次函数。

\begin{enumerate}
 \item 证明 \(n\) 次齐次函数 \(f\) 满足方程 \[
            x\frac{\partial f}{\partial x}
            +
            y\frac{\partial f}{\partial y}
            =
            nf.
          \]
 \item 利用上述性质，对于 \(z=\sqrt{x^2+y^2}\)，求 \[
            x\frac{\partial z}{\partial x}
            +
            y\frac{\partial z}{\partial y}.
          \]
 \end{enumerate}

\begin{hintbox}
对 \(f(tx,ty)=t^n f(x,y)\) 两边关于 \(t\) 求导，再令 \(t=1\)。
\end{hintbox}
\end{exercisebox}

\begin{exercisebox}
\textbf{习题 10\badge{二阶混合偏导}}

设

\[
        z=f\left(xy,\frac{x}{y}\right)+g\left(\frac{x}{y}\right),
      \]

其中 \(f\) 具有二阶连续偏导数，\(g\) 具有二阶连续导数，求 \(\frac{\partial^2 z}{\partial x\partial y}\)。

\begin{hintbox}
令 \(s=xy,\ t=x/y\)。先求 \(\frac{\partial z}{\partial y}\)，再对 \(x\) 求偏导。不要漏掉 \(f\) 的两个自变量都在变化。
\end{hintbox}
\end{exercisebox}

\begin{exercisebox}
\textbf{习题 11\badge{向量值函数复合}}

设向量值函数 \(f:R^2\to R^3\) 的坐标分量函数为

\[
        x=u^2+v^2,\qquad
        y=u^2-v^2,\qquad
        z=uv.
      \]

向量值函数 \(g:R^2\to R^2\) 的坐标分量函数为

\[
        u=r\cos\theta,\qquad v=r\sin\theta.
      \]

求复合函数 \(f\circ g\) 的导数。

\begin{hintbox}
写 \(J_f(u,v)\) 和 \(J_g(r,\theta)\)，再用 \(J_{f\circ g}=J_f(g(r,\theta))J_g(r,\theta)\)。
\end{hintbox}
\end{exercisebox}

\begin{exercisebox}
\textbf{习题 12\badge{依赖图求导}}

设 \(w=f(x,u,v),\ u=g(y,z),\ v=h(x,y)\)，求

\[
        \frac{\partial w}{\partial x},\qquad
        \frac{\partial w}{\partial y},\qquad
        \frac{\partial w}{\partial z}.
      \]

\begin{hintbox}
画依赖图：\(x\) 直接影响 \(w\)，也通过 \(v\) 影响 \(w\)；\(y\) 通过 \(u,v\) 影响 \(w\)；\(z\) 只通过 \(u\) 影响 \(w\)。
\end{hintbox}
\end{exercisebox}

\begin{exercisebox}
\textbf{习题 13\badge{极坐标型变量}}

设

\[
        z=u^v,\qquad
        u=\ln\sqrt{x^2+y^2},\qquad
        v=\arctan\frac{y}{x},
      \]

求 \(dz\)。

\begin{hintbox}
先对 \(z=u^v\) 取对数：\(\ln z=v\ln u\)。再用一阶全微分形式不变性。
\end{hintbox}
\end{exercisebox}

\begin{exercisebox}
\textbf{习题 14\badge{一阶与二阶}}

设

\[
        z=(x^2+y^2)e^{\arctan(y/x)},
      \]

求 \(dz\) 和 \(\frac{\partial^2 z}{\partial x\partial y}\)。

\begin{hintbox}
可令 \(r^2=x^2+y^2,\ \theta=\arctan(y/x)\)，先求一阶微分；二阶混合偏导再从一阶偏导继续求。
\end{hintbox}
\end{exercisebox}

\begin{exercisebox}
\textbf{习题 15\badge{全微分}}

求下列函数的全微分：

\begin{enumerate}
 \item \(u=f(ax^2+by^2+cz^2)\)；
 \item \(u=f(x+y,xy)\)；
 \item \(u=f(\ln(1+x^2+y^2+z^2),e^{x+y+z})\)。
 \end{enumerate}

\begin{hintbox}
对一阶全微分，先设中间变量再写 \(du=\sum \frac{\partial f}{\partial s_i}ds_i\) 通常最清楚。
\end{hintbox}
\end{exercisebox}

\begin{exercisebox}
\textbf{习题 16\badge{高阶全微分}}

设 \(f(t)\) 具有任意阶连续导数，而 \(u=f(ax+by+cz)\)。对任意正整数 \(k\)，求 \(d^ku\)。

\begin{hintbox}
令 \(s=ax+by+cz\)，则 \(ds=a\,dx+b\,dy+c\,dz\)，且 \(d^2s=0\)。因此 \[
          d^ku=f^{(k)}(s)(ds)^k.
        \]
\end{hintbox}
\end{exercisebox}

\begin{exercisebox}
\textbf{习题 17\badge{径向不变}}

设函数 \(z=f(x,y)\) 在全平面上有定义，具有连续的偏导数，且满足方程

\[
        x\frac{\partial f}{\partial x}(x,y)
        +
        y\frac{\partial f}{\partial y}(x,y)=0.
      \]

证明：\(f(x,y)\) 为常数。

\begin{hintbox}
固定方向 \((x,y)\)，看一元函数 \(\phi(t)=f(tx,ty)\)。由题设可得 \(\phi'(t)=0\)，所以沿每条射线常值；再用原点把所有射线连起来。
\end{hintbox}
\end{exercisebox}

\begin{exercisebox}
\textbf{习题 18\badge{Hadamard 公式}}

设 \(n\) 元函数 \(f\) 在 \(R^n\) 上具有连续偏导数，证明对于任意的

\[
        x=(x_1,\ldots,x_n),\qquad y=(y_1,\ldots,y_n)\in R^n,
      \]

成立下述 Hadamard 公式：

\[
        f(y)-f(x)
        =
        \sum_{i=1}^{n}
        \int_0^1
        (y_i-x_i)
        \frac{\partial f}{\partial x_i}(x+t(y-x))\,dt.
      \]

\begin{hintbox}
令 \(\phi(t)=f(x+t(y-x))\)。用一元微积分基本定理： \[
          f(y)-f(x)=\phi(1)-\phi(0)=\int_0^1\phi'(t)\,dt.
        \]
\end{hintbox}
\end{exercisebox}

来源：陈纪修、于崇华、金路《数学分析》第三版下册，第十二章 §2，教材第 137-146 页。下一节：第十二章 §3：中值定理和 Taylor 公式。

上一节：第十二章 §1：偏导数与全微分。复盘：错题集与好题集。路线图：教材路线图。
